% GENERATED FILE. Edit canonical lessons, then run npm run generate:books.
% canonical-source-hash: fnv1a64:ba53f6f8843d1154
% canonical-lessons: PA-C59-dhona, PA-C59-hassna, PA-C59-rona, PA-C59-gauna, PA-C59-nachchna

\chapter{To wash, To laugh, To cry, To sing, To dance}
\label{ch:pa-to-wash-to-laugh-to-cry-to-sing-to-dance}
\hlchaptermodality{59}

\begin{chapteropening}
\textbf{By the end of this chapter:} \emph{I can say to wash, to laugh, to cry, to sing and to dance.}

Complete the last lesson of chapter 59: I can say to wash, to laugh, to cry, to sing and to dance.
\end{chapteropening}

\section[\texorpdfstring{\pa{ਧੋਣਾ} (dhoṇā)}{ਧੋਣਾ (dhoṇā)}]{\pa{ਧੋਣਾ} (dhoṇā) — to wash}

\label{lesson:PA-C59-dhona}



\begin{quote}
Before the new one: say the Punjabi for to open, then the Punjabi for to cook.
\end{quote}

\subsection*{You'll want to know: \pa{ਧੋਣਾ}}
\textbf{\pa{ਧੋਣਾ}} — \emph{dhoṇā} — \textquotedblleft{}to wash\textquotedblright{}.

\begin{cousinweb}
\textbf{\pa{ਧੋਣਾ}} (\emph{dhoṇā}) is a tadbhava, a word inherited and reshaped by sound change. It comes from \emph{*dhauvati}, \textquotedblleft{}washes\textquotedblright{}, a form CDIAL reconstructs from its descendants (CDIAL 6886).
\end{cousinweb}

\begin{grammarlens}[title={What you've built}]
One more word for this chapter.
\end{grammarlens}

\subsection*{Your turn}
\begin{itemize}
\raggedright
  \item \emph{Say it:} \emph{dhoṇā}
  \item \emph{Say it:} \emph{dhoṇā}, once more
  \item \emph{Say it:} say \emph{dhoṇā}
  \item \emph{Recall:} say the Punjabi for to open, then the Punjabi for to cook, then say \emph{dhoṇā} again
\end{itemize}

\begin{tcolorbox}[breakable,colback=teal!4,colframe=teal!35!black,title={Before you move on}]
What does \pa{ਧੋਣਾ} mean? (\textquotedblleft{}To wash\textquotedblright{}.) Say it once more.
\end{tcolorbox}
\section[\texorpdfstring{\pa{ਹੱਸਣਾ} (hassṇā)}{ਹੱਸਣਾ (hassṇā)}]{\pa{ਹੱਸਣਾ} (hassṇā) — to laugh}

\label{lesson:PA-C59-hassna}



\begin{quote}
Before the new one: say the Punjabi for to cook, then the Punjabi for to wash.
\end{quote}

\subsection*{You'll want to know: \pa{ਹੱਸਣਾ}}
\textbf{\pa{ਹੱਸਣਾ}} — \emph{hassṇā} — \textquotedblleft{}to laugh\textquotedblright{}.

\begin{cousinweb}
\textbf{\pa{ਹੱਸਣਾ}} (\emph{hassṇā}) is a tadbhava, a word inherited and reshaped by sound change. It comes from \emph{*hasyate}, \textquotedblleft{}one laughs\textquotedblright{}, a form CDIAL reconstructs from its descendants (CDIAL 14048). It is built on the Sanskrit root \emph{has}, \textquotedblleft{}laugh\textquotedblright{}.
\end{cousinweb}

\begin{grammarlens}[title={What you've built}]
One more word for this chapter.
\end{grammarlens}

\subsection*{Your turn}
\begin{itemize}
\raggedright
  \item \emph{Say it:} \emph{hassṇā}
  \item \emph{Say it:} \emph{hassṇā}, once more
  \item \emph{Say it:} say \emph{hassṇā}
  \item \emph{Recall:} say the Punjabi for to cook, then the Punjabi for to wash, then say \emph{hassṇā} again
\end{itemize}

\begin{tcolorbox}[breakable,colback=teal!4,colframe=teal!35!black,title={Before you move on}]
What does \pa{ਹੱਸਣਾ} mean? (\textquotedblleft{}To laugh\textquotedblright{}.) Say it once more.
\end{tcolorbox}
\section[\texorpdfstring{\pa{ਰੋਣਾ} (roṇā)}{ਰੋਣਾ (roṇā)}]{\pa{ਰੋਣਾ} (roṇā) — to cry}

\label{lesson:PA-C59-rona}



\begin{quote}
Before the new one: say the Punjabi for to wash, then the Punjabi for to laugh.
\end{quote}

\subsection*{You'll want to know: \pa{ਰੋਣਾ}}
\textbf{\pa{ਰੋਣਾ}} — \emph{roṇā} — \textquotedblleft{}to cry\textquotedblright{}.

\begin{cousinweb}
\textbf{\pa{ਰੋਣਾ}} (\emph{roṇā}) is a tadbhava, a word inherited from Sanskrit and reshaped by sound change. It comes from \emph{roditi}, \textquotedblleft{}howls, weeps\textquotedblright{} (CDIAL 10840).
\end{cousinweb}

\begin{grammarlens}[title={What you've built}]
One more word for this chapter.
\end{grammarlens}

\subsection*{Your turn}
\begin{itemize}
\raggedright
  \item \emph{Say it:} \emph{roṇā}
  \item \emph{Say it:} \emph{roṇā}, once more
  \item \emph{Say it:} say \emph{roṇā}
  \item \emph{Recall:} say the Punjabi for to wash, then the Punjabi for to laugh, then say \emph{roṇā} again
\end{itemize}

\begin{tcolorbox}[breakable,colback=teal!4,colframe=teal!35!black,title={Before you move on}]
What does \pa{ਰੋਣਾ} mean? (\textquotedblleft{}To cry\textquotedblright{}.) Say it once more.
\end{tcolorbox}
\section[\texorpdfstring{\pa{ਗਾਉਣਾ} (gāuṇā)}{ਗਾਉਣਾ (gāuṇā)}]{\pa{ਗਾਉਣਾ} (gāuṇā) — to sing}

\label{lesson:PA-C59-gauna}



\begin{quote}
Before the new one: say the Punjabi for to laugh, then the Punjabi for to cry.
\end{quote}

\subsection*{You'll want to know: \pa{ਗਾਉਣਾ}}
\textbf{\pa{ਗਾਉਣਾ}} — \emph{gāuṇā} — \textquotedblleft{}to sing\textquotedblright{}.

\begin{cousinweb}
\textbf{\pa{ਗਾਉਣਾ}} (\emph{gāuṇā}) is inherited. CDIAL lists it under \emph{gāpayati}, the second form it gives beside Sanskrit \emph{gāyati}, \textquotedblleft{}sings\textquotedblright{} (CDIAL 4135).
\end{cousinweb}

\begin{grammarlens}[title={What you've built}]
One more word for this chapter.
\end{grammarlens}

\subsection*{Your turn}
\begin{itemize}
\raggedright
  \item \emph{Say it:} \emph{gāuṇā}
  \item \emph{Say it:} \emph{gāuṇā}, once more
  \item \emph{Say it:} say \emph{gāuṇā}
  \item \emph{Recall:} say the Punjabi for to laugh, then the Punjabi for to cry, then say \emph{gāuṇā} again
\end{itemize}

\begin{tcolorbox}[breakable,colback=teal!4,colframe=teal!35!black,title={Before you move on}]
What does \pa{ਗਾਉਣਾ} mean? (\textquotedblleft{}To sing\textquotedblright{}.) Say it once more.
\end{tcolorbox}
\section[\texorpdfstring{\pa{ਨੱਚਣਾ} (nachchṇā)}{ਨੱਚਣਾ (nachchṇā)}]{\pa{ਨੱਚਣਾ} (nachchṇā) — to dance}

\label{lesson:PA-C59-nachchna}



\begin{quote}
Before the new one: say the Punjabi for to cry, then the Punjabi for to sing.
\end{quote}

\subsection*{You'll want to know: \pa{ਨੱਚਣਾ}}
\textbf{\pa{ਨੱਚਣਾ}} — \emph{nachchṇā} — \textquotedblleft{}to dance\textquotedblright{}.

\begin{cousinweb}
\textbf{\pa{ਨੱਚਣਾ}} (\emph{nachchṇā}) is a tadbhava, a word inherited from Sanskrit and reshaped by sound change. It comes from \emph{nṛtyati}, \textquotedblleft{}dances\textquotedblright{} (CDIAL 7583).
\end{cousinweb}

\begin{grammarlens}[title={What you've built}]
One more word for this chapter.
\end{grammarlens}

\subsection*{Your turn}
\begin{itemize}
\raggedright
  \item \emph{Say it:} \emph{nachchṇā}
  \item \emph{Say it:} \emph{nachchṇā}, once more
  \item \emph{Say it:} say \emph{nachchṇā}
  \item \emph{Recall:} say the Punjabi for to cry, then the Punjabi for to sing, then say \emph{nachchṇā} again
\end{itemize}

\begin{tcolorbox}[breakable,colback=teal!4,colframe=teal!35!black,title={Before you move on}]
What does \pa{ਨੱਚਣਾ} mean? (\textquotedblleft{}To dance\textquotedblright{}.) Say it once more.
\end{tcolorbox}
